\documentclass[11pt]{article}
\usepackage[margin=1in]{geometry}
\usepackage{amsmath,amssymb}
\usepackage{booktabs}
\usepackage[dvipsnames]{xcolor}
\usepackage{hyperref}
\hypersetup{colorlinks=true,linkcolor=blue!50!black,urlcolor=blue!50!black,citecolor=blue!50!black}

\title{Counting Sources, Not Bursts:\\
Source Clustering in the Repeater/Non-Repeater Dispersion-Measure\\
Comparison of CHIME/FRB Catalog~1}
\author{David Lindgreen}
\date{October 9, 2026 \\[2pt] \small Version 0.1.0 \quad Preprint, not peer reviewed}

\emergencystretch=3em
\begin{document}
\maketitle

\begin{abstract}
We report a preregistered, independent re-analysis of one comparison in the first CHIME/FRB fast radio burst
catalog: do repeating and apparently non-repeating bursts differ in dispersion measure (DM), intrinsic pulse width
and spectral bandwidth? On 497 analysed bursts (59 bursts from 18 repeating sources, 438 non-repeaters), width and
bandwidth differ clearly (Kolmogorov--Smirnov $p<10^{-8}$), as the catalog paper reports. The DM comparison,
however, is highly significant at burst level ($p=2.0\times10^{-10}$) where the catalog paper, which uses one burst per
repeating source, reports consistent distributions. A post-hoc diagnostic explains the difference: bursts from the same
repeating source share essentially one DM excess (intraclass correlation $0.99999$), so the 59 repeater bursts behave like
about $7$ independent observations. At the independent unit, the 18 sources, the DM difference is marginal (KS $p=0.041$;
Mann--Whitney $p=0.026$) and not significant at $p<0.01$, while the width and bandwidth differences survive every
source-level test. We also correct an error in this project's earlier documentation, which wrongly attributed the
discrepancy to two prolific sources. The result is a statement about the unit of analysis, not about whether repeaters
form a distinct population.
\end{abstract}

\section{Question}
The first CHIME/FRB catalog~\cite{chime2021} lists 536 bursts detected between 2018-07-25 and 2019-07-01, 62 of them
from 18 previously known repeating sources. Its repeater/non-repeater comparisons use each repeating source's
first-detected burst, so that every source counts once. We asked a narrow question with the same public data: does
that comparison hold if every analysed burst is used instead, under a protocol fixed before any result was generated?
The question is one of statistical design, since repeat bursts from one source are not independent draws from the
population of sources~\cite{hurlbert1984}. We make no claim that repeaters and non-repeaters are, or are not,
physically distinct populations.

\section{Data and preregistration}
We used the machine-readable catalog table at CDS VizieR (\texttt{J/ApJS/257/59/table2}), selecting one row per burst.
Structural checks run at fetch time confirm 536 bursts and 62 repeater bursts from 18 sources. Following the catalog's
own flag, 39 bursts from non-nominal telescope operation were excluded, leaving 497 bursts: 59 from 18 repeating
sources and 438 non-repeaters. The protocol (committed before the registry) fixed: the DM excess after Galactic
subtraction with NE2001 as the primary measure (raw fitburst DM and YMW16 as robustness checks), intrinsic pulse width,
and bandwidth $=$ high $-$ low frequency; a two-sample Kolmogorov--Smirnov (KS) test as the primary test with
Anderson--Darling~\cite{scholz1987} and a 10{,}000-resample bootstrap of the median difference as checks. Two
hypotheses were stated: \textbf{H1}, DM distributions are indistinguishable; \textbf{H2}, width and bandwidth differ.
The registry regenerates identically from a fresh download.

\section{Preregistered results}
\begin{table}[h]
\centering\small
\begin{tabular}{lrrrrr}
\toprule
Measure & $n$ (rep./non) & KS $D$ & KS $p$ & Median rep. & Median non-rep.\\
\midrule
DM excess, NE2001 (primary) & 59/438 & 0.457 & $2.0\times10^{-10}$ & 157.0 & 510.7\\
DM, raw fitburst & 59/438 & 0.479 & $2.0\times10^{-11}$ & 349.7 & 572.8\\
DM excess, YMW16 & 59/438 & 0.487 & $8.1\times10^{-12}$ & 164.6 & 498.7\\
Pulse width (ms) & 59/438 & 0.433 & $2.4\times10^{-9}$ & 2.00 & 0.93\\
Width, excl.\ limit-flagged & 59/415 & 0.425 & $5.9\times10^{-9}$ & 2.00 & 0.97\\
Bandwidth (MHz) & 59/438 & 0.522 & $1.3\times10^{-13}$ & 206.0 & 346.1\\
Bandwidth, excl.\ limit-flagged & 59/415 & 0.529 & $6.4\times10^{-14}$ & 206.0 & 351.5\\
\bottomrule
\end{tabular}
\caption{Preregistered burst-level comparisons (DM in pc\,cm$^{-3}$). Anderson--Darling $p\le0.001$ in every row.}
\end{table}
\textbf{H2 is confirmed}: repeaters have wider pulses and narrower bandwidth, robust to every variant. \textbf{H1 is
falsified} at burst level, under all three DM conventions, with repeaters at lower DM (median DM excess 157.0 against
510.7\,pc\,cm$^{-3}$; bootstrap 95\% interval of the difference 167.4--389.4).

\section{Post-hoc diagnostics: source clustering}
The H1 result differs from the catalog paper's, and the catalog paper counts sources, not bursts. We therefore ran, after
the primary result, four fixed diagnostics (not preregistered; all reported, for all three measures): (i) the
intraclass correlation (ICC, one-way random effects~\cite{shrout1979}) of each measure within sources, with the Kish
design effect and effective sample size~\cite{kish1965}; (ii) a source-level test in which each repeating source is
collapsed to the median of its bursts; (iii) 10{,}000 random draws of one burst per source (the catalog paper's
first-detection rule is one such draw); (iv) the burst-level KS test with each source removed in turn.

\begin{table}[h]
\centering\footnotesize\setlength{\tabcolsep}{4pt}
\begin{tabular}{lccc}
\toprule
 & DM excess (NE2001) & Pulse width & Bandwidth\\
\midrule
Within-source ICC & 0.99999 & 0.19 & 0.28\\
Design effect (Kish) & 8.36 & 2.39 & 3.07\\
Effective $n$ of 59 bursts & 7.1 & 24.7 & 19.2\\
Source-level KS $p$ (18 sources) & 0.041 & $1.3\times10^{-5}$ & $4.4\times10^{-6}$\\
Source-level Mann--Whitney $p$ & 0.026 & $4.7\times10^{-6}$ & $1.4\times10^{-5}$\\
1 burst/source, KS $p$: median (max) & 0.041 (0.041) & $2.6\times10^{-4}$ (0.046) & $1.5\times10^{-5}$ ($2.5\times10^{-3}$)\\
Leave-one-source-out KS $p$ & $2\times10^{-11}$--$4.4\times10^{-5}$ & $3\times10^{-10}$--$1.6\times10^{-7}$ & $2\times10^{-14}$--$8\times10^{-11}$\\
\bottomrule
\end{tabular}
\caption{Post-hoc source-clustering diagnostics (seed 20261009). Cluster sizes: 19, 8, 3, 3, twelve sources with 2, two with 1.}
\end{table}

\paragraph{Findings.}
\emph{DM:} within a repeating source the DM excess varies by only a few pc\,cm$^{-3}$ against hundreds between sources, so the
ICC is $0.99999$ and the 59 repeater bursts carry the information of about 7 independent observations. Every one-burst-per-source
draw gives the same $p=0.041$, since each source has effectively one DM. At the independent unit, the difference is marginal:
KS $p=0.041$ and Mann--Whitney $p=0.026$, nominally significant at $\alpha=0.05$ but not at the catalog paper's $p<0.01$
convention. The median of the 18 source medians is 358.7 against 510.7\,pc\,cm$^{-3}$ for non-repeaters. The burst-level
significance therefore comes from counting each source's DM several times, with weights set by burst count, not from a few
dominant sources: removing any single source, including the most prolific (19 bursts), leaves the burst-level difference
significant ($p\le4.4\times10^{-5}$). \emph{Width and bandwidth:} ICCs are low to moderate, and every source-level test, every
random one-burst-per-source draw ($p\le0.046$ for width, $\le2.5\times10^{-3}$ for bandwidth) and every leave-one-out
test remains significant. The replicated contrast is not an artifact of clustering.

\section{Discussion}
The burst-level DM discrepancy is real in the data but is an artifact of the unit of analysis: the catalog paper's choice to
count sources is the statistically appropriate one for DM, and under it the repeater--non-repeater difference is at most
marginal in this catalog. This does not show that repeaters and non-repeaters share a DM distribution: with 18 sources the
test has little power, the source-level $p$-values sit near conventional thresholds, and a different DM convention, exclusion
or added source could move them across. A later CHIME/FRB study of 25 newly discovered repeaters reports lower mean DM
for repeaters~\cite{chime2023}; that is directionally consistent with our median difference but uses a different sample and
does not validate any burst-level $p$-value here. The accurate summary is that the DM evidence is sample- and
selection-dependent. A practical lesson follows for catalog-level comparisons: report the number of independent units,
and report the ICC when units contribute several observations.

\section{Correction to earlier project documentation}
This project's earlier report and web pages stated that two repeaters (FRB~20180916B and FRB~20180814A) supplied 33 and 20 of
the 59 repeater bursts, ``90\%'' of the sample, and called that the likely driver of the discrepancy. While preparing this paper we
recounted from the verified catalog: the counts are 19 and 8, i.e.\ 27 of 59 (46\%). The statement was wrong and has been
corrected (erratum in the repository report). No registry value, test statistic or hypothesis disposition changed.

\section{Limitations}
One catalog, one telescope and band (400--800\,MHz), one year. Repeater is a detection-window label, not a guaranteed physical
category. CHIME/FRB selection effects in DM, width and sky position are not modelled; catalog-derived columns are used as given.
Multi-component bursts are reduced to their first sub-burst. The clustering diagnostics are post-hoc, descriptive, and use small
numbers of sources; we claim an explanation of the discrepancy, not a measurement of any population difference.

\section*{Reproducibility and AI assistance}
Code, protocol, frozen registries, the clustering diagnostics and tests:\\
\url{https://github.com/lindgreendavid/frb-atlas}\\
Interactive laboratory:\\
\url{https://frb-atlas-interactive.lindgreendavid.workers.dev}\\
The analysis code, tests and drafts of this text were produced with substantial assistance from an AI system (Claude, Anthropic);
the author is responsible for the content. The protocol commit precedes the registry, and the post-hoc diagnostics are in a separate,
labelled file.

\begin{thebibliography}{9}
\bibitem{chime2021} CHIME/FRB Collaboration (M.~Amiri et al.), The first CHIME/FRB fast radio burst catalog,
\emph{Astrophys. J. Suppl.} \textbf{257}, 59 (2021), doi:10.3847/1538-4365/ac33ab; erratum \textbf{264}, 53 (2023)
(all-sky rate only), doi:10.3847/1538-4365/acb54c.
\bibitem{chime2023} CHIME/FRB Collaboration, CHIME/FRB discovery of 25 repeating fast radio burst sources,
\emph{Astrophys. J.} (2023), doi:10.3847/1538-4357/acc6c1.
\bibitem{hurlbert1984} S.~H. Hurlbert, Pseudoreplication and the design of ecological field experiments,
\emph{Ecol. Monogr.} \textbf{54}, 187--211 (1984).
\bibitem{scholz1987} F.~W. Scholz and M.~A. Stephens, $K$-sample Anderson--Darling tests,
\emph{J. Am. Stat. Assoc.} \textbf{82}, 918--924 (1987).
\bibitem{shrout1979} P.~E. Shrout and J.~L. Fleiss, Intraclass correlations: uses in assessing rater reliability,
\emph{Psychol. Bull.} \textbf{86}, 420--428 (1979).
\bibitem{kish1965} L.~Kish, \emph{Survey Sampling}, Wiley (1965).
\end{thebibliography}
\end{document}
